% Fresh quantified join of available proof bodies; not a historical-byte restoration.
\section{The complete C14.4 source union and the original C14.5 join}
\label{sec:fresh-c144-c145-join}

\paragraph{Input and scope convention.}
Let $K_c,K_t$ be arbitrarily chosen, separately labelled real closed
extensions of $\mathbb R$; no algebraic-disjointness assumption is made.
Use the original complete parent and canonical native C1 contact,
with $d>0$, $\lambda=n=0$, $m,C,A,\alpha>0$, $D,\beta\ge0$
and the original allowed literal zero convention.  Preserve physicality,
stationary contact, all strict native gaps, the positive infinitesimal
total scale, and every common entry predicate of the historical C1.4
row.  In particular $x=k^2>0$ is infinitesimal and
$\zeta>0$ is noninfinitesimal.  No source contact is identified with
an independent testing point.  Normalize $d=1$ only by the permitted
positive pure-source factor.

The exact historical row is the C1.4 row of
\path{proofs/routing/routing_audit.tex}, sections
``One exact input, before optional exclusions'', ``Exhaustive dispatch
into the historical thirteen rows'', and ``C1.4: exhaustive C14
subdivision, with missing exits restored''.

The companion
\path{ROUTING_LEDGER.json} records the common input $Q$ and each
literal source predicate.  This section completes that entered row;
it does not replace the separate proof that an arbitrary physical
support reaches the row or an earlier terminal.

At the C14 subdivision itself, the earlier exits are not left
implicit: they are the three $(L,\eta)$ admissions
$(\text{unbounded},\text{unit})$,
$(\text{unit},\text{infinitesimal})$,
$(\text{unit},\text{unit})$, the JR/JB bounded-$M_\beta$ exit,
and RO, JO, HO at unit $\eta,\zeta$.  Their exact predicates and
proof bodies are displayed below.  Any additional upstream
prior-complement predicate is carried as the same conjunct in
every alternative.  No inequality used here is inferred merely
from being outside an earlier admission.

\subsection{C14.4: an exhaustive source-only partition}
In the original C14.4 part, $L=x\zeta$ and $\sigma=1/\eta$ are
positive infinitesimals.  Set
\[
 S=\sigma+L,\qquad \kappa=x/S,\qquad
 \nu=\sigma/S,\qquad\ell=L/S,
 \qquad \nu+\ell=1,\quad\zeta=\ell/\kappa.          \tag{JC1}
\]
The exact source is FM2--FM3, obtained algebraically from NC1--NC4;
the relevant gap and sphere estimates are NC5 and ES1--ES10/SE1--SE9.
These are actual-source statements, including $D=0$ or $\beta=0$
where allowed.  The NC5 bound $E=1-\kappa h/\nu\ge1/18$ follows
from native signs and $D\ge0$ with $k\le1$, not from a prior-exit
exclusion.  Thus its use introduces no additional source predicate.

The parameter $\kappa$ is bounded.  In fact with
$M=1+xr-\kappa(1-h)/\ell>0$ and $Q_s=1-\ell Sr\sim1$,
the exact identity is
\[
 \kappa Q_s+\ell M+\nu E=1,
 \qquad 0<\kappa Q_s<1-\nu/18.                    \tag{JC2}
\]
Also $0<\nu,\ell<1$.  Each of these positive bounded source
quantities is either infinitesimal or a positive bounded unit.
The two fractions cannot both be infinitesimal.  If $\kappa$ is a
unit, then $\ell$ is a unit because $\ell=\kappa\zeta$ and
$\zeta$ is noninfinitesimal; consequently $\zeta$ is also a unit.
In particular the case ``unit $\kappa$, infinitesimal $\ell$'' is
empty, rather than omitted.

The following five alternatives are exhaustive on JC1:
\[
\begin{array}{c|c|c|l}
 \kappa&\nu&\ell&\text{complete full-parent theorem}\\ \hline
 \text{unit}&\text{unit}&\text{unit}&\mathrm{FU}\quad
       \text{(with UKR for its matched chart)}\\
 \text{unit}&\text{infinitesimal}&\text{unit}&\mathrm{NU}\\
 \text{infinitesimal}&\text{unit}&\text{unit}&\mathrm{FM}\\
 \text{infinitesimal}&\text{unit}&\text{infinitesimal}&\mathrm{EL}\\
 \text{infinitesimal}&\text{infinitesimal}&\text{unit}
       &\mathrm{SN},\mathrm{BN},\mathrm{UN}\text{ as in JC4}.
\end{array}                                                     \tag{JC3}
\]
Every entry is an actual theorem on all independent tests and all
states, not only its scalar contact graph.  FU and NU retain the
possibly vanishing $M$ and the whole divergent $m$ packet.
FM, SN, BN, UN and EL use only the valid small-$e$ source, since
SE4 proves $e\in\operatorname{Inf}$ exactly when
$\kappa\in\operatorname{Inf}$.

To spell out the last line, use only in this line
\[
 q_\nu=\nu/\ell,\qquad \tau=h/q_\nu,
 \qquad n_s=\beta/q_\nu^2.                         \tag{JC4}
\]
If $n_s$ is bounded, SN applies and its source identity gives
$\tau\sim1$.  If $n_s$ is unbounded and $\tau$ is bounded,
BN applies.  If $\tau$ is unbounded, UN applies, including
$\theta=1-h$ infinitesimal.  These alternatives are exhaustive;
BN in fact proves the entire bounded-$\tau$ domain and includes
SN as its bounded-$n_s$ part.  There is no assumption that
$h/q_\nu$ or $\beta/q_\nu^2$ remains a unit.  The EL proof uses
the reciprocal source ratio $q_\ell=\ell/\nu$ in its own formulas;
no formula is transferred between these two local definitions.

The literal faces are preserved.  On FM and EL, $D/S$ has positive
unit residue, so $D=0$ has empty intersection.  On SN,
$D/(Sq_\nu^2)\sim1/2$, with the same conclusion.  FU, NU, BN and
UN retain their actual sphere alternatives and include $D=0$
where it occurs.  Literal $\beta=0$ belongs to SN in the
double-infinitesimal endpoint, since then $n_s=0$ and the exact
boundary relation makes $\tau$ bounded.  It has empty intersection
with UN's unbounded-$\tau$ source; all other relevant theorems
retain it explicitly.  No positive opening of $r,1-r,\theta,M$
or either sphere coordinate has been removed.

Each component proof includes all testing charts.  In the small-$e$
sources, the actual SE10--SE12 homogeneous parent covers Bob residues
unequal to $\bar\ell$; CE covers noninfinitesimal Alice tests.
The remaining charts are proved by the displayed complete normal
allocations.  FU and NU supply their own raw finite/reciprocal
testing joins on the unit-$\kappa$ source.  Thus JC3--JC4 cover
every original C14.4 source/test/state tuple, including literal
axes, projective infinities, all angle signs and independent
real and imaginary state components.  This proves the whole
original C14.4 branch.

\subsection{C14.1 and the retained complete C14.2/C14.3 branches}
For $L,\eta$ unbounded set $\sigma=1/\eta$ and $H=h/\sigma$.
The actual large-$L$ source formulas are those of P1/SP6 and SF1--SF5.
If $H$ is bounded, the surviving
\path{c141a_closure_and_c141b_partial.tex} theorem proves
all tests and states.  Its proof uses JR/JB when
$M_\beta=\beta(1-h)/\sigma^2$ is bounded, including zero,
and its own full state-free normal conductor for the remaining
unbounded-$M_\beta$ sector.  Thus the original C14.1a branch is
entirely covered without reinstating a positive-$\beta$ premise.

For unbounded $H$, the surviving SF theorem first removes
$\operatorname{res}h<1/2$.  On the other source part, US bounds
$k/\sigma$ and covers its positive-unit case, equivalently the
unit beta-sphere coordinate.  The bounded-state, bounded testing
ratio, unit-Alice, literal and sign exits are the actual BS, BC,
US, GA and inherited CE proofs indexed in
\path{proofs/c1_unbounded/DEPENDENCIES.txt}.  Their retained
complement is exactly contained in
\[
 \operatorname{res}h\ge1/2,\quad k/\sigma\in\operatorname{Inf},
 \quad b_0\in\operatorname{Inf},\quad d_0\sim1,
 \quad u,w>0\text{ infinitesimal},                 \tag{JC5}
\]
with the original positive-state and unbounded-ratio predicates.
The fresh AR1--AR21 theorem covers this complement using its restored
angular SOS budget.  Its coefficients require no bound on
$qu/w$, $u/w$, $W/Z$ or $uW/(wZ)$, and all SP6 first-normal
terms are allocated.  Applying the real routing to each state
component independently closes the whole original C14.1b branch.
Its literal $D=0$ source uses the US sphere endpoint; its
$\beta=0$ source is empty when $H$ is unbounded by the exact sphere
boundary equation.  These statements keep the original source
intersection rather than asserting a missing face is positive.

The complete C14.2 proof remains the actual R19
\path{c142_all_tests_join.tex}, with its bounded-$H$,
sphere, toric, state-ratio and LA/LB leaves.  It covers the original
$L$-unit, $\eta$-unbounded domain, both literal coefficient faces,
all openings and all independent tests/states.  The actual R18
\path{c143ab_all_tests_join.tex}, together with its completed
$j$ conductor and source/remote leaves, covers C14.3a and C14.3b.
These complete proof bodies are present; no missing late-R21
theorem statement is substituted for them.

\subsection{The original C1.4 source is exhausted}
Put $L=x\zeta$.  The exact NI4--NI6 source implications and the
explicit C14 subdivision in the current routing proof give the
following elementary decision tree.

If $L$ is unbounded, $Q=1-Lr>0$ implies $r$ infinitesimal.
The native identity $1-r=a_{\rm nat}\eta$ with
$0<a_{\rm nat}<2$ therefore excludes infinitesimal $\eta$.
Unit $\eta$ is the complete
\path{advances/c1_large_l_eta_unit/large_l_unit.tex}
theorem.  Unbounded $\eta$ is split by bounded or unbounded $H$,
just covered by C14.1a/JR/JB and C14.1b.

If $L$ is a positive bounded unit, the three possible types of
$\eta$ give respectively the complete
\path{advances/c1_lunit/lunit_eta_small.tex} theorem,
\path{advances/c1_lunit_eta_unit/eta_unit.tex} theorem,
and C14.2.  These statements retain the actual $D,\beta$ hypotheses
and their full remote joins.

If $L$ is infinitesimal, unbounded $\eta$ is the now complete
C14.4 branch.  For bounded $\eta$, first split $\zeta$:
unbounded $\zeta$ is C14.3a; unit $\zeta$ with infinitesimal
$\eta$ is C14.3b.  The only remaining source has both
$\eta,\zeta$ positive bounded units.  Its exact positive gap is
\[
 j=\eta(1+xh)-1+r=\eta J>0,
 \qquad j_0=\eta+r-1,
 \qquad j-j_0=\eta xh\in\operatorname{Inf}.        \tag{JC6}
\]
The gap $j$ is bounded.  If it is a unit, $j_0$ is a positive
unit and RO applies.  If $j$ is infinitesimal, then
$\eta+r\sim1$ and $1-r$ is a unit; unit $r$ gives JO,
infinitesimal $r$ gives HO.  Their actual bodies are
\path{advances/revision13_c1_small_l/}
\path{r_opening_closure.tex},
\path{j_opening_closure.tex}, and
\path{simultaneous_r_j_closure.tex}.
They retain actual $j>0$, even when $j_0=0$ or slightly negative,
and compare neither $j$ with $x$ nor a source quantity with a
testing quantity.  RO contains the earlier IC1 regular-gap
admission, so that historical fourth C1.4 subdomain is restored
as part of a fully stated theorem.

These alternatives exhaust every positive $L,\eta,\zeta$ permitted
in the entered C1.4 row.  More formally, let $\mathcal C$ denote
that exact common row predicate and let $E_{14}$ be the disjunction
of the seven explicitly named prior exits above.  Let $R_i$ be the
six original residual predicates after their stated exclusions.
The unchanged finite source argument proves
\[
 \mathcal C\ \Longrightarrow\ E_{14}\ \lor\
 R_{14.1a}\lor R_{14.1b}\lor R_{14.2}\lor
 R_{14.3a}\lor R_{14.3b}\lor R_{14.4}.              \tag{JC7}
\]
Each alternative now has an actual complete full-parent proof.
Intersecting all terms with any fixed earlier entry predicate
preserves this implication; the predicate is not dropped or used
to manufacture a source bound.  Conversely, every prior C14 exit
removed from an $R_i$ is restored by its displayed proof.  This is
the required quantified C14.5 union, including all source openings,
allowed literal coefficient zeros, projective endpoints and exact
sign gauges.  Multiplying by the original positive pure-source
normalization and positive pure-testing homogenization factors
preserves finite separated SOS.

\paragraph{Checklist consequence and boundary of this result.}
C14.1, C14.4 and C14.5 are complete on their unchanged original
criteria.  C14.2 and C14.3 retain their existing complete status.
Consequently U1.B1, the whole original entered C1.4 row, is complete.
This join does not by itself prove a global physical-source entry,
the independent weak-C2 row, the finite-lift bridge, or an evaluated
affine LMI.  It adds no uniform square-count hypothesis.
The accompanying dependency receipt binds the actual old and new
proof bodies.  Focused symbolic source identities and finite
decision-tree checks corroborate the explicit analytical union;
neither historical completion claims nor assertion counts are
mathematical premises.
